\documentclass[11pt, a4paper]{article}

\usepackage[T1]{fontenc}
\usepackage[utf8]{inputenc}
\usepackage{tgpagella}
\usepackage[spanish, es-noshorthands]{babel}
\usepackage{amsmath, amssymb, bm}
\usepackage{tabularx}
\usepackage{booktabs}
\usepackage[most]{tcolorbox}
\usepackage[margin=2.5cm]{geometry}
\usepackage{ragged2e}
\usepackage{fancyhdr}
\usepackage[american]{circuitikz}

\pagestyle{fancy}
\fancyhf{}
\setlength{\headheight}{16pt}
\lhead{\textbf{UCB Sede Tarija} | Ing. Mecatrónica}
\chead{}
\rhead{IMT-121 Circuitos Electrónicos I | Examen EC2}
\lfoot{Prof: M.Sc. Ing. Bernardo Quiroga Turdera}
\rfoot{Página \thepage}
\renewcommand{\headrulewidth}{0.4pt}

\definecolor{ucbblue}{RGB}{0, 51, 102}

\begin{document}

\begin{tcolorbox}[colback=ucbblue!5!white, colframe=ucbblue, title=\textbf{Evaluación Escrita Individual: Hito 2 (Unidad EC2)}]
    \centering \textbf{Teoremas de Redes y Simplificación Estática}
\end{tcolorbox}

\vspace{0.2cm}
\textbf{Estudiante:} \hrulefill \quad \textbf{Fecha:} \hrulefill

\vspace{0.2cm}
\noindent\textit{Instrucciones: Duración 45--60 min. No se permite calculadora ni formulario. La evaluación prioriza el método, la desactivación correcta de fuentes y el manejo de referencias/signos por sobre la aritmética. Muestre su procedimiento completo. (Se asume la relación de potencia: $P_L = V_L \cdot I_L$ y $P_L(R_L) = V_{th}^2 \frac{R_L}{(R_{th}+R_L)^2}$).}

\vspace{0.3cm}
% =====================================
% PARTE A
% =====================================
\section*{Parte A: Superposición y Referencias (25 pts)}
Para la siguiente red lineal, se solicita determinar la tensión del nodo $a$ con respecto a tierra ($V_a$) utilizando el Teorema de Superposición.

\begin{center}
\begin{circuitikz}[scale=1, transform shape]
    % Fuente izquierda y R1 protegidos con llaves {...}
    \draw (0,0) to[V, l={$12\,\text{V}$}, invert] (0,2) to[R, l={$R_1=2\,\text{k}\Omega$}] (3,2);
    % Rama central a tierra
    \draw (3,2) to[short, *-o] (3,2) node[above]{$a$};
    \draw (3,2) to[R, l={$R_3=2\,\text{k}\Omega$}] (3,0);
    % Fuente derecha y R2. Protegidos con llaves {...}
    \draw (6,0) to[V, l_={$-6\,\text{V}$}, invert] (6,2) to[R, l_={$R_2=2\,\text{k}\Omega$}] (3,2);
    
    % Tierra
    \draw (0,0) -- (6,0);
    \draw (3,0) node[ground]{};
    
    % Flecha indicativa de Va
    \draw[->, thick, blue] (2.5, 1.8) -- (2.5, 0.2) node[midway, left]{$V_a$};
\end{circuitikz}
\end{center}

\begin{enumerate}
    \item Defina y documente la desactivación de la fuente de $-6\,\text{V}$. Calcule la contribución de la fuente de $12\,\text{V}$ ($V_a^{(1)}$).
    \item Defina y documente la desactivación de la fuente de $+12\,\text{V}$. Calcule la contribución de la fuente de $-6\,\text{V}$ ($V_a^{(2)}$).
    \item Obtenga el valor total de $V_a$.
    \item En una sola oración, interprete físicamente qué significa que una de las contribuciones calculadas haya resultado negativa respecto a la referencia solicitada.
\end{enumerate}

% =====================================
% PARTE B
% =====================================
\section*{Parte B: Equivalencias de Puerto (Thévenin y Norton) (35 pts)}
Considere el siguiente circuito de alimentación con una resistencia de carga extraíble $R_L$ conectada en el puerto $a-b$ (donde $b$ es la referencia de tierra).

\begin{center}
\begin{circuitikz}[scale=1, transform shape]
    % Todo protegido con llaves {...}
    \draw (0,0) to[V, l={$V_s=12\,\text{V}$}, invert] (0,3) to[R, l={$R_1=4\,\text{k}\Omega$}] (4,3);
    \draw (4,3) to[R, l={$R_2=4\,\text{k}\Omega$}] (4,0);
    \draw (4,3) to[short, *-o] (6,3) node[above]{$a$};
    \draw (0,0) -- (6,0) node[below]{$b$} to[short, *-o] (6,0);
    \draw (2,0) node[ground]{};
    \draw[dashed, blue] (6,3) -- (7.5,3) to[R, l={$R_L$}] (7.5,0) -- (6,0);
\end{circuitikz}
\end{center}

\begin{enumerate}
    \item Retire la carga $R_L$ conceptualmente y determine la tensión de circuito abierto ($V_{th}$).
    \item Desactive correctamente las fuentes independientes y determine la resistencia equivalente mirando desde el puerto $a-b$ ($R_{th}$). 
    \item Dibuje explícitamente el equivalente de Thévenin.
    \item Determine la corriente de corto circuito de puerto ($I_N$) y dibuje el equivalente de Norton.
\end{enumerate}

% =====================================
% PARTE C
% =====================================
\section*{Parte C: Efecto de Carga y Máxima Transferencia (40 pts)}
A partir de este punto, asuma un Equivalente de Thévenin caracterizado por: $\mathbf{V_{th} = 12\,\text{V}}$ y $\mathbf{R_{th} = 4\,\text{k}\Omega}$. 

\subsection*{C1. Evaluación Analítica Limpia}
Se conectan sucesivamente tres resistencias de carga $R_L$: \textbf{$2\,\text{k}\Omega$, $4\,\text{k}\Omega$ y $8\,\text{k}\Omega$}.
\begin{enumerate}
    \item Calcule la corriente de carga ($I_L$), la tensión de carga ($V_L$) y la potencia consumida ($P_L$) para \textbf{cada uno} de los tres resistores. Muestre su procedimiento.
    \item A partir de sus cálculos, encierre en un círculo cuál carga produce el máximo voltaje, cuál la máxima corriente y cuál la máxima potencia. ¿Por qué estas tres condiciones no ocurren con el mismo resistor?
\end{enumerate}

\subsection*{C2. Interpretación Experimental (Datos Simulados de Evaluación)}
La siguiente tabla presenta valores obtenidos por un grupo de estudiantes en el laboratorio midiendo una red cuyo $R_{th}$ teórico y medido fue \textbf{$4\,\text{k}\Omega$}.

\vspace{0.2cm}
\noindent
\begin{tabularx}{\textwidth}{*{4}{>{\centering\arraybackslash}X}}
\toprule
\textbf{$R_L$ (k$\Omega$)} & \textbf{$V_L$ medido (V)} & \textbf{$I_L$ medido (mA)} & \textbf{$P_L$ medido (mW)} \\
\midrule
$2$ & $3.9$ & $1.95$ & $7.6$ \\ \addlinespace
$4$ & $5.8$ & $1.45$ & $8.4$ \\ \addlinespace
$5$ & $6.6$ & $1.32$ & $8.7$ \\ \addlinespace
$8$ & $7.8$ & $0.98$ & $7.6$ \\
\bottomrule
\end{tabularx}
\vspace{0.2cm}

\begin{enumerate}
    \item ¿En qué resistencia $R_L$ ocurre el máximo experimental de potencia? Compare este hallazgo con la predicción teórica ($R_{th} = 4\,\text{k}\Omega$).
    \item ¿Esta discrepancia observada invalida el modelo matemático de la máxima transferencia? Proponga \textbf{una hipótesis técnica} que justifique físicamente el desplazamiento del máximo y diga cómo lo comprobaría experimentalmente.
\end{enumerate}

\end{document}
